\section{Finite sets of primes}
\label{sec:support-prime-no-go}

The local columns are indexed by primes.  This section records the local
analogue of the previous one: a fixed finite set of primes cannot certify
a statement about all primes.  It then describes the numerical support
table used for the examples, and says precisely what that table does and
does not establish.

\begin{theorem}[A finite set of primes does not decide all-prime coverage]\label{thm:apparatus:finite-prime-cover-no-go}
Let a \emph{coverage record} assign to each prime $p$ a status, covered or
uncovered, and let the all-prime target be the statement that every prime
is covered.  For every finite set $P_0$ of primes there is a prime
$q\notin P_0$, and there are two coverage records that agree at every
prime of $P_0$, such that the first covers every prime and the second
leaves $q$ uncovered.  Consequently no function of the restriction of a
record to $P_0$ decides the all-prime target.%
\footnote{Lean proves the statement for coverage records indexed by natural
numbers, with an explicit prime $q$ dividing $1+\prod_{p\in P_0}p$, and
with primality decided by a finite divisor check.  These are formal
records; no claim is made that they arise from elliptic curves.}
\end{theorem}

\begin{proof}
Let $q$ be a prime factor of $1+\prod_{p\in P_0}p$; then $q\notin P_0$.
Let the first record cover every prime and let the second agree with it
except at $q$.  The two records have the same restriction to $P_0$, so
any function of that restriction gives them the same value, while the
all-prime target holds for the first and fails for the second.
\end{proof}

For computations one works with a finite truncation of the local column.

\begin{definition}[Support-prime truncation residual]\label{def:apparatus:support-prime-truncation-residual}
Fix a curve $E$ and a set $R(E)$ of primes, the \emph{residual rows}: the
primes at which no local comparison is taken as available.  For $B\ge1$
the truncation residual is
\[
  \Xi_B(E)=\#\{p\le B:\ p\text{ prime},\ p\in R(E)\}.
\]
It depends only on $R(E)$ below $B$.  In the examples below $R(E)$ is the
set of primes of bad reduction.  For the five window curves and $B=49$
this gives
\[
  \Xi_{49}(11a1)=\Xi_{49}(37a1)=\Xi_{49}(121b1)=1,\qquad
  \Xi_{49}(960d1)=3,\qquad
  \Xi_{49}(571a1)=0,
\]
the last becoming $1$ once $B\ge571$.  This is a finite diagnostic,
distinct from \Cref{thm:apparatus:finite-prime-cover-no-go}.
\end{definition}

Taking $R(E)$ to be the bad primes treats every good prime as covered by
an ordinary or signed local theory.  For good ordinary primes this is the
setting of \Cref{thm:apparatus:ordinary-schur-collapse}.  For good
supersingular primes, the signed theory of
\Cref{def:apparatus:p-adic-map} requires $p$ odd and $a_p=0$, and the
table below shows that this fails in some small cases.

\begin{remark}[Numerical support table]\label{rmk:apparatus:support-prime-atlas}
For the five window curves and the $15$ primes below $50$, together with
the pair $(571a1,571)$, we have computed a $76$-row table of
$(E,p)$ data from the Cremona models: the number of points of the
reduction and $a_p$ (each obtained by two independent methods), the
reduction type (good, multiplicative or additive), and for good primes
whether the reduction is ordinary or supersingular.  The table is
included with the supporting files of the paper.  The bad primes are
\[
  11a1:\{11\},\quad
  37a1:\{37\},\quad
  121b1:\{11\},\quad
  960d1:\{2,3,5\},\quad
  571a1:\{571\}.
\]
Among the good supersingular rows, five lie outside the basic numerical
scope ($p$ odd and $a_p=0$) of Kobayashi's signed
theory~\citep{Kobayashi2003}: $(11a1,2)$ and $(37a1,2)$ with $a_2=-2$;
$(37a1,3)$ with $a_3=-3$; and $(121b1,2)$ and $(571a1,2)$ with $a_2=0$.
For these rows the cited signed input does not supply a local comparison.
Generalized signed constructions that allow $p=2$ or $a_p\neq0$ exist,
but would have to be specified and checked row by row.  If these five
rows are counted as residual, the truncation residuals below $50$ become
$2,3,2,3,1$ for $11a1,37a1,121b1,960d1,571a1$.  The table certifies
the numerical local data only; it does not certify any local comparison
or main-conjecture input for any row.
\end{remark}

The finiteness statements of this paper are collected in
\Cref{tab:no-go-suite}.
\begin{table}[tbp]
\centering
\footnotesize
\caption{The apparatus no-go suite.  Each item is an audit result about what
the displayed public row does not determine; none is a non-existence theorem
about richer external source data.}
\label{tab:no-go-suite}
\begin{tabularx}{\textwidth}{@{}>{\raggedright\arraybackslash}p{0.15\textwidth}
>{\raggedright\arraybackslash}p{0.27\textwidth}
>{\raggedright\arraybackslash}p{0.36\textwidth}
>{\raggedright\arraybackslash}X@{}}
\toprule
No-go & Form & What the literature or public row lacks & Status \\
\midrule
NG-1 & Finite scalar shadow & The scalar finite factor does not identify
which finite-source block carries the valuation. & mechanized finite
noninjectivity witness. \\
\addlinespace
NG-2 & \(\dim \Sha[p]\) shadow & Dimension-only data forgets the
Cassels--Tate determinant and pairing structure. & mechanized projector
trace residual. \\
\addlinespace
NG-3 & Gram-matrix/regulator shadow & A regulator determinant does not
construct the rank-\(2\) source whose Gram matrix it would measure. &
mechanized typed shadow statement. \\
\addlinespace
NG-4 & Finite-window promotion & A bounded audit window cannot certify an
unbounded operator family. & mechanized finite-window obstruction. \\
\addlinespace
NG-5 & Finite-prime coverage & A fixed finite prime set cannot determine an
all-prime support carrier. & mechanized finite-prime cover obstruction. \\
\addlinespace
Higher-rank GZ & Rank-\(\geq2\) source gap & The theorem-grade
generalized Gross--Zagier identity is not supplied by the audited inputs. &
recorded as an audit-result gap. \\
\addlinespace
Global audit & Operator promotion gap & Local or finite scalar evidence does
not by itself promote to the global determinant-line carrier. & recorded as
scope boundary. \\
\addlinespace
Support-prime atlas & Local support gap & The \(b50\) support table is finite
evidence, not all-prime coverage. & support-only audit table. \\
\bottomrule
\end{tabularx}
\end{table}

